\section{Bayesian benchmarks for the frequency-format experiments}
\label{sec:annex-p19-natural-frequencies}

\begin{annexitem}{annex:p19.frequency-bayes-bridge}{thm:p19-frequency-bayes-bridge}{The joint-law bridge and normalization}
\begin{proof}
For \(0<\pi<1\), both hypotheses have positive probability, so the
likelihoods are defined. Apply \cref{thm:p15-finite-bayes} to the partition
\(H,H^c\) and evidence \(D\). Its two intersection masses are
\(x=\pi u\) and \(z=(1-\pi)v\); total probability gives \(s=x+z\).
Since \(s>0\), the posterior is \(x/s\). Substitution yields the
standard-menu formula, and the short and hybrid formulas are the same
identity with the supplied \(x\) or the product \(\pi u\).
If a menu independently supplies \(s\), equivalence requires that it equal
\(\pi u+(1-\pi)v\); there is no permission to disregard a discrepancy.

Take a collection of \(N\) distinguishable cases and give each mass
\(1/N\). Partition them into the four cells of sizes \(a,b,c,d\).
Finite additivity gives \(\mathbb P(H\cap D)=a/N\) and
\(\mathbb P(D)=(a+c)/N\). When \(a+c>0\), the definition of
conditional probability in \cref{thm:p15-finite-conditioning-laws} gives
\[
 \mathbb P(H\mid D)=\frac{a/N}{(a+c)/N}=\frac{a}{a+c}.
\]
This direct calculation also covers empty hypothesis rows. If \(H\)
is empty, \(a=0\) and the posterior is zero; if \(H^c\) is empty,
\(c=0\) and the posterior is one. A likelihood conditional on an empty
row is not introduced. If \(a+c=0\), the denominator is zero and this
calculation makes no conditional-probability assertion.

When both rows are nonempty, their rates satisfy
\[
 \pi u=\frac{a+b}{N}\frac{a}{a+b}=\frac aN,\qquad
 (1-\pi)v=\frac{c+d}{N}\frac{c}{c+d}=\frac cN.
\]
This proves the bridge to the two-hypothesis Bayes formula explicitly.
For rational \(\pi,u,v\), all four numbers
\[
 \pi u,\quad\pi(1-u),\quad(1-\pi)v,\quad(1-\pi)(1-v)
\]
are rational, nonnegative and sum to one. Choose a positive integer \(N\)
divisible by their denominators and multiply them by \(N\) to obtain
\(a,b,c,d\). Their row sums are \(N\pi\) and \(N(1-\pi)\), which
are positive, and the preceding rate calculation recovers the triple.
Conversely, every proportion in a finite uniform collection is rational,
which explains the restriction on an exact integer representation.

For a positive integer \(k\),
\(ka/(ka+kc)=a/(a+c)\), proving invariance under common scaling.
If instead both rows are normalized to a common size, their positive
fractions are proportional to \(u\) and \(v\); pooling those synthetic
rows gives \(u/(u+v)\). With \(u,v>0\) the denominators are positive,
and cross multiplication gives
\[
 \frac{\pi u}{\pi u+(1-\pi)v}=\frac{u}{u+v}
 \quad\Longleftrightarrow\quad
 \pi uv=(1-\pi)uv
 \quad\Longleftrightarrow\quad \pi=\tfrac12.
\]
Thus the equal-row calculation need not preserve the intended prior.
All these identities concern a stipulated or empirical joint law; they
do not identify the distribution of a new case without additional premises.
\end{proof}
\end{annexitem}

\begin{annexitem}{annex:p19.bayesian-shortcut-errors}{thm:p19-bayesian-shortcut-errors}{Errors, signs and explicit small-ratio bounds}
\begin{proof}
With positive displayed denominators, subtraction of the rare-event
approximation from the exact posterior gives
\[
 \frac{x}{x+z}-\frac{x}{x+v}
 =\frac{x(v-z)}{(x+z)(x+v)}
 =\frac{x\pi v}{(x+z)(x+v)},
\]
using \(z=(1-\pi)v\). This is nonnegative. For the high-hit-rate
approximation, use \(x=\pi u\) to obtain
\[
 \frac{\pi}{\pi+z}-\frac{x}{x+z}
 =\frac{z(\pi-x)}{(\pi+z)(x+z)}
 =\frac{\pi z(1-u)}{(\pi+z)(x+z)}\geq0.
\]
For the small-\(x\) comparison approximation,
\[
 \frac{x}{z}-\frac{x}{x+z}=\frac{x^2}{z(x+z)}.
\]
For the small-\(z\) version,
\[
 \frac{x}{x+z}-\left(1-\frac{z}{x}\right)
 =\frac{x^2-(x-z)(x+z)}{x(x+z)}
 =\frac{z^2}{x(x+z)}.
\]
Setting \(r=x/z\) in the first error expresses it as
\(r^2/(1+r)\). For \(0\leq r\leq\varepsilon\), it is at most
\(r^2\leq\varepsilon^2\). Setting \(r=z/x\) proves the second
bound in exactly the same algebraic form. The divisions require \(z>0\)
or \(x>0\), respectively, as stated; zero denominators are not limits
silently substituted for defined values.

The combined formulas mentioned in the body are obtained by replacing
\(1-\pi\) by one and \(u\) by one in the relevant comparison
approximation. No error bound for them is asserted merely from qualitative
words such as rare or high. For example, with \(\pi=.8,u=.5,v=.2\),
the formula \(\pi/v\) is 4, despite the exact posterior being
\(.4/(.4+.04)=10/11\). With \(\pi=.1,u=.5,v=.8\), the reverse
formula \(1-v/\pi\) is \(-7\), whereas the posterior is
\(.05/(.05+.72)=5/77\). Both triples have positive evidence probability;
the approximation's failure is not a null-event issue.
\end{proof}
\end{annexitem}

\begin{annexitem}{annex:p19.frequency-numerical-reconstructions}{thm:p19-frequency-numerical-reconstructions}{All task ratios and the two measles versions}
\begin{proof}
For the mammography standard-probability inputs,
\(\pi u=.01\times.8=.008\) and
\((1-\pi)v=.99\times.096=.09504\). Their sum is .10304 and
\(.008/.10304=800/10304=25/322\). The exact scaled count table sums
to \(100+25+1188+11187=12500\). Its hypothesis row has 125 cases,
of which 100 are positive; its other row has 12375, of which 1188 are
positive. Hence \(125/12500=.01\), \(100/125=.8\), and
\(1188/12375=.096\). The count posterior is \(100/(100+1188)=25/322\).

The source's smaller table sums to \(8+2+95+895=1000\), with 103
positive cases and 8 positive \(H\)-cases, giving \(8/103\). The short
probability formula is \(.008/.103=8/103\), so it agrees exactly with
that table. The false-positive rate \(95/990\) differs from .096 by
\(1/24750\), making the intended rounding explicit. Equal-sized
synthetic rows give \(800/(800+96)=25/28\), as claimed.

For general supplied pairs \(h/N,a/h,c/m\), multiply the prior and
hit rate to get \(a/N\). The other joint mass is
\((N-h)c/(Nm)\). Their ratio therefore is
\[
 \frac{a/N}{a/N+(N-h)c/(Nm)}
 =\frac{am}{am+(N-h)c}.
\]
The hypotheses on the rates and positive evidence ensure that this division
is defined. Treating \(a,c\) instead as two joint-cell counts gives
\(a/(a+c)\). With \(m=N-h>0\), the common factor \(m\) cancels
and the formulas coincide.

Here are all fifteen substitutions from source Table~2. The middle column
treats the printed hit and false-positive numerators as joint leaves; the
right column instead uses every printed denominator literally as a rate.
The five marked rows have a rounded \(H^c\) denominator. These two columns
are computed quantities, distinct from the source's rounded percentage
column in the body.
\begin{center}\small
\begin{tabular}{lrr}
Problem&Joint-leaf calculation&Literal-rate calculation\\\hline
Mammography&\(8/103\)&\(8/103\)\\
Measles, Study~1\(^{*}\)&\(1/6\)&\(2000/11979\)\\
Cab&\(12/29\)&\(12/29\)\\
HIV\(^{*}\)&\(1/11\)&\(1000/10999\)\\
Needle mark\(^{*}\)&\(1/20\)&\(10000/199981\)\\
Pregnancy&\(19/24\)&\(19/24\)\\
Driver&\(11/111\)&\(11/111\)\\
Posture&\(2/21\)&\(2/21\)\\
School accident&\(27/415\)&\(27/415\)\\
Suicide\(^{*}\)&\(3/10003\)&\(15/50003\)\\
Ball&\(12/13\)&\(12/13\)\\
Economics&\(3/8\)&\(3/8\)\\
Feminist&\(1/101\)&\(1/101\)\\
Rolex\(^{*}\)&\(2/27\)&\(1600/21599\)\\
Admission&\(135/199\)&\(135/199\)
\end{tabular}
\end{center}
For example, the Study~1 measles rates give
\(10\cdot10000/[10\cdot10000+9979\cdot50]=2000/11979\).
Multiplication by 100 yields approximately 16.695884, which rounds to
the source's 16.70\%; the leaf value gives 16.666667\%. The other rows
follow by the displayed substitutions, without assuming that an intentionally
rounded row total equals the complement of its printed hypothesis count.

The Study~2 measles numbers are different. They give
\[
 x=.005\times.4=.002,\qquad
 z=.995\times.0001=.0000995,\qquad s=.0020995.
\]
Consequently \(\theta=.002/.0020995=4000/4199\), approximately
.95260776. The printed hybrid menu has \(\pi=.005,u=.4,s=.0021\),
and the short menu has \(x=.002,s=.0021\); both therefore yield
\(20/21\), approximately .95238095. The difference originates in
rounding the evidence probability, not in a failure of Bayes' formula.

An exact count realization of the standard menu takes \(N=2000000\)
with \((a,b,c,d)=(4000,6000,199,1989801)\). The row totals are
10000 and 1990000; their hit rates are .4 and .0001, and their positive
total is 4199. A short-menu realization can instead take a collection
of 100000 cases with 210 positives, of which 200 are in \(H\). For
example, \((a,b,c,d)=(200,300,10,99490)\) also matches the hybrid
\(\pi=.005,u=.4,s=.0021\). Its false-positive rate is
\(10/99500\), not exactly .0001.

The rare-event shortcut for the standard inputs is
\(.002/(.002+.0001)=20/21\). The high-hit-rate shortcut is
\(.005/(.005+.0000995)=10000/10199\), approximately .98048828.
The reverse-comparison shortcut is
\(1-.0000995/.002=.95025\); the combined reverse shortcut is
\(1-.0001/.005=.98\). These computations explain why the rare-event
shortcut meets the source's 95\%/96\% rounding allowance, while a 98\%
answer does not. Their accuracy is evaluated for the actual task numbers.

For the distinct HIV shortcut example, \(x=.002\times.99=.00198\)
and \(z=.998\times.02=.01996\). Therefore
\(\theta=.00198/.02194=99/1097\), approximately .09024613, whereas
\(.002/.02=.10\). For mammography the combined shortcut is
\(.01/.096=5/48\), approximately .10416667. Its exact benchmark
\(25/322\) was derived above. These task versions remain distinct from
the HIV row of Table~2 and from the rounded frequency mammography table.

Finally, the Appendix's complement calculation satisfies
\(10/210=1/21\), or approximately 4.761905\%; its complement is
\(20/21\), or approximately 95.238095\%. The source's 95.7\% remains
the recorded answer, not the corrected division. Multiplying the three
hybrid decimal probabilities gives
\(.005\times.0021\times.4=.0000042\), or .00042\%.
The short product \(.0021\times.002\) is the same. Thus the recorded
.04\% in the two multiply-all cases is not the correct product under
the stated units. Correcting these benchmark computations does not alter
the source's historical response records or its coding rule.
\end{proof}
\end{annexitem}
